\homework{6}{Tue 11 Oct 2022 22:28}{}

\begin{enumerate}
	\item 16.A. Let $x_1 \in \boldsymbol{R}$ satisfy $x_1>1$ and let $x_{n+1}=2-1 / x_n$ for $n \in \boldsymbol{N}$. Show that the sequence $\left(x_n\right)$ is monotone and bounded. What is its limit?
\begin{proof}
	We establish the fact that $(x_n)$ is monotone decreasing and bounded below by $1$. The first term $x_1$ is obviously monotone decreasing (since there is only one term), and bounded below by $1$ as assumed.

	Now assume that $x_1, \ldots, x_k$ is monotone decreasing and bounded below by $1$. We have
	\begin{align*}
		&x_{k+1}< x_k \\
		&\iff 2 - \frac{1}{x_k} < x_k \\
		&\iff 2 x_k - 1 < x_k ^2 &\text{since }x_k > 1 > 0 \\
		&\iff  (x_k - 1)^2 > 0
	.\end{align*}
	Which is true since $x_k > 1$, thus $x_k - 1>0$.

	Moreover, since $x_k > 1$, $\frac{1}{x_k}< 1$, so $2 - \frac{1}{x_k} > 1$. Therefore $x_{k+1} > 1$. 

	We have completed the induction and $(x_n)$ is indeed monotone decreasing and bounded below by $1$. To obtain the limit, let the limit be $x$ and take limit on both sides of $x_{n+1}=2-1 / x_n$. Hence $x = 2 - 1 /x$. Solve the equation we obtain $x = 1$, so the limit is $1$.
\end{proof}

	\item 16.E. Show that every sequence in $\boldsymbol{R}$ either has a monotone increasing subsequence or a monotone decreasing subsequence.
\begin{proof}
	(Simpler)
	Take any sequence $X$. If $X$ is unbounded, then we may assume without loss of generality that it is unbounded from above, thus we can construct a monotone increasing subsequence of $X$. For the detailed construction, let
	\[
		n_j = \operatorname{inf} 
		\{n \in \mathbb{N}: x_{n}>x_{n_{j-1}}, n > n_{j-1}\} 
	.\] 
	We verify that $n_j$ is well-defined. Since $X$ is unbounded from above, for any fixed $x_{n_{j-1}}$, $x_n > x_{n_{j-1}}$ for sufficiently large $n$. Since such $n$ is infinitely many, we can add the restriction $n > n_{j_-1}$. Then by well-orderedness of natural numbers, $n_j$ is well defined.
	Then $x_{n_j}$ is monotone increasing by construction.

	If $X$ is bounded, then by Bolzano-Weierstrass, there exists convergent subsequence $(x_{n_m})$ of $X$ with limit $x$. We define $A_1 = \{m \in \mathbb{N}: x_{n_m} < x\} $, $A_2 = \{m \in \mathbb{N}: x_{n_m} = x\} $, and $A_3 = \{m \in \mathbb{N}: x_{n_m} > x\} $. Since $A_1 \cup A_2 \cup A_3 = \mathbb{N} $ is infinite, at least one of $A_1, A_2, A_3$ is infinite. If $A_2$ is infinite, we can choose a constant subsequence of $x_{n_m}$ which is indeed monotone.

	The cases $A_1$ is infinite and $A_3$ infinite are symmetrical, so it suffices to prove for the case that $A_1$ is infinite. Then we define the subsequence $x_{n_{m_j}}$ as
	\[
		m_j = \operatorname{inf} 
		\{m \in \mathbb{N}: m > m_{j-1}, x_{n_{m_{j-1}}}< x_{n_m} < x\} 
	.\] 
	To justify that $m_j$ is well-defined, note that since  $A_1$ is infinite, there exists infinitely many $x_{n_m} < x$. Since  $x_{n_m} \to x$ as $m \to \infty $, for any fixed $x_{n_{m_{j-1}}}$, we have $x_{n_m}>x_{n_{m_{j-1}}}$ for all sufficiently large $m$. Since there are infinitely many such $m$, we can add the restriction that $m > m_{j-1}$. Then by well-orderedness of natural numbers, $m_j$ is well-defined.

	Now $x_{n_{m_j}}$ is strictly monotone increasing by construction.
\end{proof}

\begin{proof}
	If $\limsup X$ is finite, then there exists a monotone subsequence $X'$ of $X$ converging to $s :=\limsup X$ (claim 1). If $\limsup X$ is not finite, then $X$ is unbounded from above, so there exists a monotone increasing unbounded subsequence of $X$ (claim 2).

	Justification of claim 1:
	Since $s = \limsup X$, there exists a subsequence of $X$ converging to $s$, say $Y$. Note that one of the following events must happen infinitely often: 
	\begin{itemize}
		\item  $y_n > s$
		\item $y_n = s$ 
		\item $y_n < s$
	\end{itemize}
	If $y_n = s$ infinitely often, we just take the subsequence that takes a constant value of $s$ and it is indeed monotone. If $y_n > s$ infinitely often, we define the subsequence $y_{n_j} $ as \[
		n_j := \operatorname{inf} \{m \in \mathbb{N}: m > n_j, s < y_m \le  y_{n_j}\} 
	.\] 
	To check the definition is well-posed, assume inductively that $n_j$ is well-defined up to $j = p$. Now since $y_n$ converges to $s$, there exists some $N \in \mathbb{N}$ such that for all  $m \ge  N$, we have $y_m \le y_{n_j} $. Since $y_m > s$ infinitely often, $\{m \in \mathbb{N}: m > n_j, s < y_m \le  y_{n_p}\}$ is nonempty, therefore (by well-orderedness of $\mathbb{N}$ ) $n_j$ is well-defined. Now $y_{n_j} $ is the desired sequence.

	The proof for the case that $y_n < s$ infinitely often is entirely symmetrical. 

	Justification of claim 2: Since $X$ is unbounded from above, for each $M$, there exists $N$ such that for all $n > N$, we have  $x_n > M$. Therefore, we can define a subsequence $x_{n_j}$ such that \[
		n_j := \operatorname{inf} 
		\{m \in \mathbb{N}: m > n_{j-1}, x_m > x_{n_{j-1}}\} 
	.\] 
	To verify that the definition is well-posed, assume inductively that $n_j$ is well-defined up to $j = p$. Now there exists $N$ such that for all $m >  N$, we have $x_m > x_{n_{p}}$. Then $\max(n_{j-1}, N)+1 \in \{m \in \mathbb{N}: m > n_{j-1}, x_m > x_{n_{j-1}}\}$. Thus by well-orderedness of natural numbers $n_{j}$ is well-defined for $j = p+1$.

	By definition, $x_{n_j}$ is monotone increasing and is a subsequence of $X$.
\end{proof}

	\item 16.G. Determine the convergence or divergence of the sequence $\left(x_n\right)$, where
\[
x_n=\frac{1}{n+1}+\frac{1}{n+2}+\cdots+\frac{1}{2 n} \quad \text { for } \quad n \in N
\]
\begin{proof}
	We establish an upper bound of the sequence:
	\[
		 x_n = \sum_{k=n+1}^{2n} \frac{1}{k}
		 \le  \sum_{k=n+1}^{2n} \frac{1}{n}
		 = n \frac{1}{n} = 1
	.\] 
	Then we prove $x_n$ is monotone increasing. Combine with the upper bound and we know $x_n$ is convergent by monotone convergence theorem.
	\begin{align*}
		x_n - x_{n-1 }&= \frac{1}{n+1} + \ldots + \frac{1}{2n} - \left( \frac{1}{n} + \ldots + \frac{1}{2n-2} \right)  \\
			      &= \frac{1}{2n-1} + \frac{1}{2n} - \frac{1}{n} \\
			      &> \frac{1}{2n} + \frac{1}{2n} - \frac{1}{n} \\
			      &= 0
	.\end{align*} 
\end{proof}

	\item 16.Q. Prove that if $K_1$ and $K_2$ are compact subsets of $\boldsymbol{R}^{\mathrm{p}}$, then there exist points $x_1 \in K_1, x_2 \in K_2$ such that if $z_1 \in K_1, z_2 \in K_2$, then $\left\|z_1-z_2\right\| \geq\left\|x_1-x_2\right\|$
\begin{proof}
	Define the distance function
		\begin{align*}
			d: K_1 \times K_2 &\longrightarrow \mathbb{R} \\
			k_1, k_2 &\longmapsto d(k_1, k_2) = \|k_1 - k_2\|
		.\end{align*}
	By definition we know $d(k_1, k_2) \ge 0$, so $d(K_1 \times K_2)$ is bounded below by $0$. Also, it is obviously nonempty. Hence there exists $u := \operatorname{inf} d(K_1\times K_2)$.

	Now for each $n \in \mathbb{N}$, there must exist $k_1^{(n)}$ and $k_2^{(n)}$ such that $d(k_1^{(n)}, k_2^{(n)}) < u + \frac{1}{n}$, by property of infimum. 

	Then since $K_1$ and $K_2$ are compact, there exist convergent subsequences $k_1^{(n_k)} \to k_1$ and $k_2^{(n_k)}\to k_2$. Thus $k_1 \in K_1$ and $k_2 \in K_2$, and $d(k_1, k_2) \le  u$. Since $d(k_1, k_2) \ge u$, we have $d(k_1, k_2) = u$. Thus for any $z_1 \in K_1, z_2 \in K_2$, we have $d(z_1, z_2) \ge d(k_1, k_2)$.
\end{proof}

\begin{proof}
	(First attempt, more complicated)
	If $K_1$ and $K_2$ intersect somewhere, then we pick $x_1 = x_2 \in K_1 \cap K_2$ and complete the proof.

	From now on we assume $K_1$ and $K_2$ are disjoint. Since $K_2$ is compact, $K_2$ is closed.
	For each $k_1 \in K_1$, we choose $k_2(k_1) \in K_2$ using Nearest Point Theorem (11.6), such that for all $k \in K_2$, $\|k_1, k\| \ge \|k_1, k_2(k_1)\|$.

	Now define \begin{align*}
		\varphi: K_1 &\longrightarrow \mathbb{R} \\
		k_1 &\longmapsto \varphi(k_1) = \|k_1 - k_2(k_1)\|
	.\end{align*}
	Let $u = \operatorname{inf}_{k_1 \in K_1}\varphi(k_1)$. To see that $u$ is well-defined, note that $u$ is nonempty (pick arbitrary $k_1 \in K_1$ and $\varphi(k_1)$ is well-defined), and bounded from below (obviously  $\varphi(k_1)\ge 0$ ). 

	Now for each $n \in \mathbb{N}$, define
	\[
		F_n = \{k_1 \in K_1: \varphi(k_1) \le  u + \frac{1}{n}\} 
	.\] 
	Then $F_n$ is closed (since it is the intersection of $K_1$ and the closed ball $\overline{B_{u+\frac{1}{n}}\left( k_2(k_1) \right) } $\quad ) and nonempty (by property of infimum). Moreover, $F_n$ has a nesting property. To see this, note that $\varphi(k_1)$ is independent of $n$, and for any $k_1 \in F_{n+1}$ we have $\varphi(k_1)\le u + \frac{1}{n+1} \le u + \frac{1}{n}$, thus $F_{n+1} \subset F_n$.

	Finally, $F_1$ is bounded since $K_1$ is bounded.
	Now by the cantor intersection theorem, there exists a point $k^*$ belonging to all the sets $F_n$. Thus $\varphi(k^*) \le u + \frac{1}{n}$ for all $n \in \mathbb{N}$, so $\varphi(k^*) = u$. We can easily verify from the definition that $ x_1:=k^*$ and $x_2 := k_2(k^*)$ satisfy desired properties.
\end{proof}

	\item 18.B. If $X=\left(x_n\right)$ is a bounded sequence in $\boldsymbol{R}$, show that there is a subsequence of $X$ which converges to $\lim \inf X$.
\begin{proof}
	Denote $x = \lim \inf X$. Let $L$ be the set of $v \in \mathbb{R}$ such that there exists a subsequence of $X$ which converges to $v$, and we know that $x = \operatorname{inf} L$. 

	Now let $\varepsilon_m = \frac{1}{m}$. Using property of infimum, take $v_m \in L$ such that $v_m < v + \varepsilon_m / 2$. Then there exists some subsequnce $(x_{n_j})$ of $X$ that converges to $v_m$. Then there exists some $N_m \in \mathbb{N}$ such that for all $j > N_m$ we have $|x_{n_j} - v_m | < \varepsilon_n / 2$.
	Combining the results, we have $|x_{n_j} - v| \le |x_{n_j} - v_m| + |v_m - v| = \varepsilon_m$. 

	Now we proceed to define our subsequence.
	For each $m \in \mathbb{N}$, take $\varphi(m)$ to be the least integer $n$ such that $n = n_j$ for some $j > N_m$, and $n > \varphi(m-1)$.

	Now we see that for each $m$, $|x_{\varphi(m)} - v|\le \varepsilon_m$. As $m \to \infty$, $\varepsilon_m \to 0$, so $x_{\varphi(m)} \to v$.
	
\end{proof}

	\item 18.F. If $X=\left(x_n\right)$ is a bounded sequence of strictly positive elements in $\boldsymbol{R}$ show that $\lim \sup \left(x_n^{1 / n}\right) \leq \lim \sup \left(x_{n+1} / x_n\right)$.
\begin{proof}
	\begin{notation}
		We say some statement $P_n$ holds for \textit{sufficiently large }$n$ if there exists $N \in \mathbb{N}$ such that for all  $n \ge  N$, we have $P_n$ is true.
	\end{notation}

	Denote $\lim \sup \frac{x_{n+1}}{x_n} = L$. Then for each $\varepsilon>0$, $\frac{x_{n+1}}{x_n}\le L + \varepsilon$ for sufficiently large $n$.
	To see this, apply the property of limsup, and see that there are at most finitely many $n$ such that $\frac{x_{n+1}}{x_n}>L + \varepsilon$.

	Thus $x_n\le   C (L + \varepsilon)^n$ for sufficiently large $n$ for some $C$ (by using a past homework problem). Therefore
	\[
		x_n ^{\frac{1}{n}} \le  C^{\frac{1}{n}}(L + \varepsilon) \to  L + \varepsilon
	.\] 
	Thus $x_n^{\frac{1}{n}}\le L + 2\varepsilon$ for sufficeintly large $n$. To see this, pick arbitrarily small $\delta>0$, and note that $C^{\frac{1}{n}}< 1 + \delta$ for sufficiently large $n$. Hence $x_n ^{\frac{1}{n}}\le (1 + \delta)(L + \varepsilon)$ for arbitrarily small $\delta$ and $\varepsilon$. Now let $\delta < \frac{1}{(L + 1)}$ and $\varepsilon< 1$, and we obtain $x_n^{\frac{1}{n}}\le L + 2 \varepsilon$ for sufficiently large $n$.

	Therefore $\lim\sup x_n^{\frac{1}{n}}\le L + 2 \varepsilon$ for arbitrary $\varepsilon$. Hence $\lim\sup x_n^{\frac{1}{n}} \le L$.
\end{proof}
\item 18.I. Show that lim $\sup X=+\infty$ if and only if there is a subsequence $X^{\prime}$ of $X$ such that $\lim X^{\prime}=+\infty$.
\begin{proof}
	Suppose $\lim \sup X = + \infty $. 
	Then for each natural number $N$, there exists some value $v > N$ in the extended real numbers ($v \in \mathbb{R}$ or $v = +\infty $), such that there is a subsequence in $X$ converging to  $v$. Thus 
	$X$ is not bounded from above, because we can always find a subsequence converging to $v$, satisfying $x_{n_j}>N$ for sufficiently large $j$.
	Then we let  \[
		n_j = \operatorname{inf} 
		\{m: m > n_{j-1}, x_m > j\} 
	.\] 
	The set above is clearly nonempty for each $j$ since $X$ is unbounded. Thus $n_j$ is well-defined and we have $x_{n_j} \to  +\infty $.

	Now assume there is a subsequence  $X'$ of $X$ such that $\lim X' = + \infty $.
	Let $V$ be the set of values $v \in \mathbb{R}$ such that there exists a subsequence of  $X$ converging to $v$. By definition, $\lim \sup X  = \operatorname{sup} V$. But we have $+\infty  \in V$, so $\lim \sup X = +\infty $.
\end{proof}

\end{enumerate}
